\section*{Présentation de l'étudiant}
\phantomsection
\addcontentsline{toc}{section}{Présentation de l'étudiant}

Je m'appelle Thomas Graber et je suis étudiant en mastère informatique, spécialité développement, programmation et data/IA, au sein de Learn IT, sur le campus Brest Open Campus. Avant cette formation, j'ai suivi une licence d'informatique à l'Université de Franche-Comté, à Besançon, où j'ai acquis des bases généralistes en systèmes, réseaux, algorithmique, bases de données et développement logiciel.

Initialement attiré par le développement web front-end, j'ai progressivement développé un intérêt plus marqué pour le développement back-end, l'architecture logicielle et la manipulation de données. Depuis le 16 septembre 2024, j'effectue mon alternance au sein du SHOM, Service hydrographique et océanographique de la Marine, en tant que développeur informatique spécialisé dans les services web.

Cette expérience professionnelle m'a permis d'évoluer dans un environnement technique complexe, sécurisé et fortement orienté métier. Elle s'inscrit dans la continuité de mon parcours de formation, au cours duquel j'ai également réalisé plusieurs projets académiques mobilisant des compétences en conception logicielle, développement, bases de données, gestion de projet et analyse de données.

Ce dossier professionnel présente donc un ensemble de missions et de projets issus de mon alternance, de ma formation en mastère et, lorsque cela est pertinent, de mon parcours antérieur en licence informatique. Ces éléments sont mobilisés au regard du référentiel RNCP du titre Expert en Architecture et Développement Logiciel, afin de démontrer les compétences acquises dans des contextes professionnels et pédagogiques complémentaires.

L'objectif de ce dossier n'est pas uniquement de décrire les travaux réalisés, mais de montrer ma capacité à analyser des problématiques logicielles réelles ou reconstituées, à concevoir des solutions adaptées, à justifier des choix techniques et à prendre du recul sur les résultats obtenus.

\clearpage
\section*{Organisation du dossier}
\phantomsection
\addcontentsline{toc}{section}{Organisation du dossier}

Afin de faire apparaître clairement les compétences mobilisées, ce dossier est structuré par blocs de compétences plutôt que par projet. Ce choix permet de relier les différentes missions et réalisations de mon parcours aux attendus de la certification.

Chaque bloc s'appuie sur des éléments de preuve issus de plusieurs contextes : missions réalisées en entreprise au SHOM, projets menés dans le cadre du mastère Learn IT, et certains travaux significatifs réalisés durant ma licence informatique. Cette organisation permet d'éviter une présentation uniquement chronologique et de montrer comment des expériences complémentaires contribuent à la validation des compétences visées.